\begin{lemma}[Situation 7.2, Subcase 3.2]
Every \noderef{HierarchySituationTwoCaseThreeTwoData} is impossible.
Equivalently, normalized Case 3 cannot occur when the last-edge color
\(\theta=\varphi_n(e_p)\) is missing at \(y_{p-1}\).
\end{lemma}
\begin{proof}
Write \(H=C.\mathit{hierarchy}\), let \(q\) be its last stored index, and
put
\[
 A=\noderef{HierarchyStarredStage}(H,q),\qquad
 \ell=\noderef{HierarchyFullTailLength}(C),
\]
\[
 B=\noderef{HierarchyFullTailPrefix}(C,\ell-2),\qquad
 B^+=\noderef{HierarchyFullTailPrefix}(C,\ell-1),\qquad
 D_{n,q}=\noderef{HierarchyStoredColorsAt}(H,q).              \tag{1}
\]
The inherited common-tail fields and \(\ell\ge2\) identify exact
\noderef{TreeSequenceInitialSegment}s
\[
 A\preccurlyeq B\prec B^+\prec T,\qquad
 |B|=|T|-2,\quad |B^+|=|T|-1.                               \tag{2}
\]
The final vertices of \(B^+\) and \(T\) are \(y_{p-1}\) and \(y_p\),
and the final individual edge is \(e_p\).  Apply
\noderef{HierarchyCounterexamplePenultimateElementary} to the six fields
of the supplied \noderef{HierarchyMinimumCounterexampleContext}.  It
gives the derived fact (6.6)
\[
             V(B^+)\text{ is elementary under }\varphi_n.    \tag{3}
\]

We record the complete normalized condition (20), which is part of the
datum rather than reproved here.  There is a vertex of \(B\) missing
\(\alpha\), \(\alpha\) is missing also at \(y_p\), and
\[
       \alpha\notin\varphi_n\langle T-A\rangle .              \tag{4}
\]
Moreover exactly one of the following certified alternatives holds:
\[
\begin{array}{ll}
q>0:&\alpha\notin D_{n,q}\cup\Gamma^q,\\
q=0:&\alpha\notin D_n\cup\Gamma^0,\\
&\alpha\in D_{n,q}\text{ and }
  \Gamma^q_\alpha\cap\varphi_n\langle T-A\rangle=\varnothing.
\end{array}                                                   \tag{5}
\]
Here \(D_n=\noderef{TashkinovStoredColors}(\mathcal E,n)\), and the last
line is the stored alternative.  This is precisely the three-way
`normalizedAlternative' field of
\noderef{HierarchySituationTwoCaseThreeNormalizedData}; in particular,
the divider-zero line concerns raw \(D_n\).

The Subcase 3.2 field says that \(\theta\) is missing at
\(y_{p-1}\).  Equation (3) shows that no vertex of \(B\), hence no vertex
of \(A\), misses \(\theta\).  Thus its
\noderef{HierarchyFirstMissingPosition} is exactly \(y_{p-1}\).  Also
\(\alpha\ne\theta\): the individual edge \(e_p\), incident with \(y_p\),
has color \(\theta\), whereas \(\alpha\) is missing at \(y_p\).

We will use the following replacement bookkeeping in every recoloring.
For an affected pair \((c,d)\), first apply \noderef{ChangeLemma} using
(3), the six context fields, the exact first-missing order and unused-edge
interval, and the applicable positive- or zero-divider alternative.  Its
component equality, \noderef{TwoColorComponentSpanningPath},
\noderef{TwoColorIncidenceDegreeAtMostTwo}, and
\noderef{TwoColorComponentsDisjoint} separate the component rooted at
\(y_p\) from the component through the two first-missing witnesses.  If
the witnesses coincide, that vertex misses both colors and is an isolated
vertex of \noderef{TwoColorIncidenceGraph}; otherwise equality would put
three distinct degree-at-most-one missing endpoints into a finite path.
This handles coincident witnesses rather than silently assuming them
distinct.  The spanning-path theorem supplies the possibly singleton
\noderef{TwoColorPathComponent} required by the stable-change providers.

When the exact Cases 1/2 proposition holds for the first color \(c\),
invoke \noderef{StableChangeCasesOneTwo} before using any stability
consequence.  If both affected colors are absent from every exact edge of
\(T-A\), except possibly the individual divider-zero SE/RE seed protected
by the returned \noderef{StableColoring}, then every tail edge is fixed.
The same component-separation argument at vertices gives equality of
missing sets on \(B^+\); unaffected stored colors retain their first
positions, and an affected stored color retains its position once its
witness lies in \(B^+\).  These are exactly the hypotheses of
\noderef{StableChangeCasesOneTwoFullTail}, which returns a
\noderef{CorrespondingExtendedTashkinovTree} on the same \(T\), maximum
property, a replacement base, a full hierarchy, a
\noderef{GoodHierarchy}, and \noderef{SameGoodHierarchy} data.

If only the prefix \(B^+\) is fixed, inspect the actual final
\noderef{TashkinovSeriesStep}.  The PE constructor gives a
\noderef{TashkinovClosure} and augmentation route from \(T_n\); the SE/RE
constructors give the exact singleton connecting edge, literal
\noderef{TreeSequenceAdjoinEdge} seed, and route from that seed.  Restrict
the old route at \(B^+\) by
\noderef{TashkinovAugmentationSequenceRestrict}.  Below \(A\), transport
it using the returned \noderef{HierarchyWeaklyStable}: for positive
divider count this is \noderef{TreeColoringInvariant}; at divider zero,
the PE case uses \noderef{CombinedStableClosurePreserved} and the exact
entry decomposition in \noderef{OrderedTreeUnion}, while the SE/RE case
has \(A=T_n\) and retains its protected exact seed.  Above \(A\), fixed
edge values and vertexwise missing-set equality allow
\noderef{TashkinovAugmentationSequenceColoringTransport}.  The single old
\noderef{TashkinovAugmentation} \(B^+\to T\), together with both affected
colors missing on \(B^+\) after the swap, supplies the suffix input.
After checking stored first positions and the affected-Gamma cutoff,
\noderef{StableChangeCasesOneTwoSplitTail} returns the same full package
without claiming that the last edge was fixed.

In either construction, the swap makes \(c\) or \(d\), as specified
below, missing at \(y_p\), and component separation preserves its earlier
witness.  These distinct witnesses give non-elementarity and hence an
actual \noderef{HierarchyCounterexample}.  Construct the new minimum
context field by field: \noderef{TwoColorComponentSwapProper} gives
properness; the pointwise swap fixes `none' and hence
\noderef{IsOneEdgeUncolored}; rung agreement transports the universally
quantified lower-rung six-pack; same-good divider equality transports
fewer-divider induction; literal terminal-tree equality and same stages
identify both coordinates of
\noderef{HierarchyCounterexampleMeasure} and transport lexicographic
minimality; and the positive terminal stability slot identifies missing
sets on \(T_n\) and transports the \(2n+11\) estimate.  This fills every
field of \noderef{HierarchyReplacementPresentation}.

For a second recoloring, the provider is initially relative to the first
presentation.  We always compose it back to the original ETT: rung,
ladder-tree, raw stored-set, connecting-edge, extension-kind, and
generating-coloring equalities compose fieldwise; positive-slot stability
composes by \noderef{sc2} after rewriting its protected stored set; and
the two \noderef{SameGoodHierarchy} relations compose their stages,
divider positions, computed stored sets, and every Gamma pair.  The final
presentation and minimum context are therefore original-relative, not
merely relative to the intermediate coloring.

We now follow the storedness cases of the manuscript.

\textbf{First suppose \(\theta\in D_{n,q}\).}
The two-color cardinality field of \(C.\mathit{good}\) supplies distinct
colors \(\gamma_1,\gamma_2\) with
\[
                 \Gamma^q_\theta=\{\gamma_1,\gamma_2\}.      \tag{6}
\]
Since \(\theta\) is not missing on \(A\) and its first missing vertex is
\(y_{p-1}\), goodness and (6) make both Gamma colors missing on the last
ordinary stage and exclude them from every exact edge of \(B^+-A\).
Neither equals \(\theta\), by the disjointness of stored colors and Gamma
colors in \noderef{GoodHierarchyReservedDisjoint}; hence neither colors
\(e_p\).  The normalized alternatives (4)--(5) also exclude \(\alpha\)
from the tail and from \(\Gamma^q_\theta\).  We have proved manuscript
equation (26):
\[
 \{\alpha,\gamma_1,\gamma_2\}\cap
       \varphi_n\langle T-A\rangle=\varnothing .              \tag{7}
\]

Apply \noderef{ChangeLemma} with ordered colors
\((\gamma_1,\alpha)\).  The first color is stage-missing; its first
witness lies in \(A\), the first alpha witness lies in \(B\), and (7)
gives the exact unused interval.  Thus the exact Cases 1/2 proposition
holds, at \(T_{n,q}\) for \(q>0\) and at \(T_n\) for \(q=0\).
Separate the component through these witnesses from the one rooted at
\(y_p\) as above, and set
\[
 \mu_1=\noderef{TwoColorComponentSwap}
       (G,\varphi_n,\gamma_1,\alpha,\operatorname{inl}y_p).   \tag{8}
\]
Apply the Cases 1/2 theorem and then the full-tail provider.  Equation (7)
fixes every exact tail edge, so there is no unverified full-tail equality;
vertexwise equality on \(B^+\) follows from component separation, and
stored first positions are unchanged because \(\gamma_1\) is stage-missing,
while a stored \(\alpha\), if present, has a witness already in \(B\).
Obtain an actual replacement presentation \(P_1\) by the bookkeeping
above.  The swap makes \(\gamma_1\) missing at \(y_p\) and preserves its
stage witness.  Moreover
\[
 \{\alpha,\gamma_1,\gamma_2\}\cap
       \mu_1\langle T-A\rangle=\varnothing,                  \tag{9}
\]
because neither affected color occurred on an old tail edge.

Relative to \(P_1\), apply \noderef{ChangeLemma} with ordered colors
\((\gamma_1,\theta)\).  The first color remains stage-missing, theta's
first missing vertex remains \(y_{p-1}\), and (9) is the unused interval.
The exact Cases 1/2 proposition again holds with \(\gamma_1\) in the
stage-missing slot.  The component through the two first witnesses is
disjoint from the component rooted at \(y_p\); define
\[
 \mu_2=\noderef{TwoColorComponentSwap}
       (G,\mu_1,\gamma_1,\theta,\operatorname{inl}y_p).       \tag{10}
\]
Invoke the Cases 1/2 theorem immediately.  This swap fixes precisely the
prefix \(B^+\): (9) excludes \(\gamma_1\), and
\noderef{HierarchyTailMissingColorExclusion} excludes theta up to its
first witness, except for the divider-zero exact seed, which the retained
stability fixes.  Missing sets agree on \(B^+\).  The actual old suffix is
the one-step addition of \(e_p\), colored theta and supported by theta's
missingness at \(y_{p-1}\).  Stored positions are unchanged; the only
affected Gamma color is \(\gamma_1\), whose pair owner is theta and whose
first missing position is \(|B^+|\).  The split-tail construction therefore
gives the second full replacement package.  Compose it back to the
original presentation using \(\noderef{sc2}\) and the fieldwise
same-good composition described above.

The final presentation satisfies
\[
 \theta\in\overline\mu_2(y_{p-1})\cap\overline\mu_2(y_p),
 \quad \theta\in D^{2}_{n,q},\quad
 \mu_2(e_p)=\gamma_1,\quad
 \gamma_1\notin\overline\mu_2(y_{p-1}).                      \tag{11}
\]
The first three assertions follow from the component swap, prefix
missing-set equality, stored-set agreement, and the last-edge change.  For
the last, elementarity of \(B^+\) under \(\mu_1\) puts the earlier
\(\gamma_1\)-missing witness away from \(y_{p-1}\), and the component
through \(y_{p-1}\) is disjoint from the one changed at \(y_p\).
The transported common-tail fields, (6), and (11) form an actual
\noderef{HierarchySituationTwoCaseOneData} for the final presentation.
It is impossible by \noderef{HierarchySituationTwoSubcaseOneOne}.

\textbf{Now suppose \(\theta\notin D_{n,q}\).}
Because theta is not missing on \(A\), expanding
\noderef{HierarchyStoredColorsAt} upgrades this to theta outside raw
\(D_n\).  Every Gamma color is stage-missing, so (3) and theta's first
witness at \(y_{p-1}\) also give \(\theta\notin\Gamma^q\).  Applying
\noderef{HierarchyTailMissingColorExclusion} at \(y_{p-1}\), the sole
possible exceptional seed would carry theta in raw \(D_n\), which has
just been excluded.  Therefore
\[
 \theta\notin D_{n,q}\cup\Gamma^q\ (q>0),\qquad
 \theta\notin D_n\cup\Gamma^0\ (q=0),                       \tag{12}
\]
and theta colors no exact edge of \(B^+-A\).  This is manuscript (28),
including its zero-divider content.

Split first on whether \(\alpha\in D_{n,q}\).

If \(\alpha\notin D_{n,q}\), the stored alternative in (5) is
impossible.  Hence (5) gives
\[
 \alpha\notin D_{n,q}\cup\Gamma^q\ (q>0),\qquad
 \alpha\notin D_n\cup\Gamma^0\ (q=0).                       \tag{13}
\]
This is manuscript (29).  Apply \noderef{ChangeLemma} to the ordered pair
\((\alpha,\theta)\).  Both colors are missing on \(B^+\), their first
witnesses occur respectively in \(B\) and at \(y_{p-1}\), and (4) gives
the exact alpha-unused interval.  For \(q>0\), either alpha is missing on
\(T_{n,q}=A\), or (12)--(13) put both colors outside \(D_{n,q}\); for
\(q=0\), either alpha is missing on \(T_n\), or both lie outside raw
\(D_n\).  These are precisely the alternatives of the change lemma, so
it identifies the component through the two first-missing witnesses.
Separate it from the component rooted at \(y_p\), and define
\[
 \sigma=\noderef{TwoColorComponentSwap}
       (G,\varphi_n,\alpha,\theta,\operatorname{inl}y_p).     \tag{14}
\]

At this point split on the exact proposition required by
\noderef{StableChangeCasesOneTwo}: either \(q>0\) and alpha is missing on
\(T_{n,q}\), or \(q=0\) and alpha is missing on \(T_n\), or alpha is
missing on \(A\) while both affected colors lie outside raw \(D_n\).
In this Cases 1/2 branch invoke the theorem immediately.  Equations
(4), (12) fix every individual edge from \(A\) through \(B^+\), missing
sets agree on that prefix, and the old suffix is the one-step addition of
the theta-colored \(e_p\).  Both affected colors are missing after the
swap at the disjoint first-witness component.  Neither is in a Gamma pair
by (12)--(13), and any affected stored first position is witnessed in
\(B^+\).  Thus the exact route, stored-position, and cutoff hypotheses of
\noderef{StableChangeCasesOneTwoSplitTail} hold, yielding an actual
full-tree replacement package.

In the complementary branch, we do not pass (12)--(13) off as
stage-missingness.  Failure of the exact proposition gives that alpha is
not missing on \(A\): this is immediate when \(q>0\); when \(q=0\), the
raw exclusions in (12)--(13) make the third disjunct true whenever alpha
is missing on \(A\), so failure rules that out as well.  Hence every
vertex of \(A\) precedes alpha's first missing witness, that witness lies
in \(B\), and it precedes theta's first witness \(y_{p-1}\).  Equations
(12)--(13) put both colors outside \(D_{n,q}\), and (4) says alpha is
unused on every exact edge between those two witnesses.  These are exactly
the hypotheses of \noderef{StableChangeCaseThree}, with ordered colors
\((\alpha,\theta)\).  That theorem directly supplies weak stability,
mod-coloring, a corresponding same-rung maximum ETT on the full exact
\(T\), a replacement base, and a same-good hierarchy.  No full- or
split-tail reconstruction is imposed in this external-first-missing
branch.

In either subbranch, build the counterexample and every field of its
minimum context by the replacement bookkeeping above.  The resulting
presentation \(P\) has
\[
 \theta\in\overline\sigma(y_{p-1})\cap\overline\sigma(y_p),
 \quad \theta\notin D^P_{n,q},\quad
 \sigma(e_p)=\alpha,\quad
 \alpha\notin\overline\sigma(y_{p-1}).                       \tag{15}
\]
The last assertion follows from elementarity (3): alpha already has a
missing witness in \(B\), and the component changed at \(y_p\) is
disjoint from the first-witness component containing \(y_{p-1}\).
Stored-set agreement transports theta's exclusion.  Thus (15) and the
transported common-tail fields form an actual
\noderef{HierarchySituationTwoCaseTwoData};
\noderef{HierarchySituationTwoSubcaseTwoOne} gives a contradiction.

It remains to consider \(\alpha\in D_{n,q}\).  Choose the two distinct
colors \(\varepsilon_1,\varepsilon_2\) in
\(\Gamma^q_\alpha\).  The first two alternatives in (5) are now
impossible, and the third gives manuscript (30):
\[
 \{\alpha,\varepsilon_1,\varepsilon_2\}\cap
       \varphi_n\langle T-A\rangle=\varnothing .              \tag{16}
\]
Apply \noderef{ChangeLemma} with ordered colors
\((\varepsilon_1,\alpha)\).  The first is stage-missing, alpha's first
witness lies in \(B\), and (16) is the unused interval; hence the exact
Cases 1/2 proposition holds.  Separate components and define
\[
 \mu_1=\noderef{TwoColorComponentSwap}
       (G,\varphi_n,\varepsilon_1,\alpha,
        \operatorname{inl}y_p).                              \tag{17}
\]
Invoke the Cases 1/2 theorem and the full-tail provider.  Equation (16)
fixes every tail edge; component separation gives missing-set equality on
\(B^+\); epsilon is not stored, and alpha's stored first position is
already witnessed in \(B\).  Hence we obtain an actual presentation
\(P_1\) with all six minimum-context assertions transported.  The swap
makes \(\varepsilon_1\) missing at \(y_p\) and preserves its stage
witness.  Since neither old affected color occurred on the tail,
\[
        \varepsilon_1\notin\mu_1\langle T-A\rangle.           \tag{18}
\]

Now apply \noderef{ChangeLemma} relative to \(P_1\), with ordered colors
\((\varepsilon_1,\theta)\).  Epsilon remains stage-missing; theta's first
witness remains \(y_{p-1}\) by prefix missing-set equality; and (18)
supplies the unused interval.  The exact Cases 1/2 proposition again holds
with epsilon in the stage-missing slot.  Separate components, put
\[
 \mu_2=\noderef{TwoColorComponentSwap}
       (G,\mu_1,\varepsilon_1,\theta,
        \operatorname{inl}y_p),                              \tag{19}
\]
and invoke the Cases 1/2 theorem before protecting the divider-zero seed.
Equations (12), (18) fix every exact edge through \(B^+\), and the same
internal-vertex argument preserves all missing sets there.  Restrict and
transport the actual PE or exact-seed SE/RE route to \(B^+\) as described
above; the old suffix is the one-step addition of theta-colored \(e_p\).
Stored first positions are unchanged, and pairwise disjointness says that
the only Gamma pair containing \(\varepsilon_1\) belongs to alpha, whose
first witness lies in \(B\), so the Gamma cutoff holds.  The split-tail
provider gives the second actual full-tree package.  Compose it back to
the original ETT and hierarchy field by field, using \(\noderef{sc2}\) at
every positive stability slot.

The final original-relative presentation satisfies
\[
 \theta\in\overline\mu_2(y_{p-1})\cap\overline\mu_2(y_p),
 \quad \theta\notin D^2_{n,q},\quad
 \mu_2(e_p)=\varepsilon_1,\quad
 \varepsilon_1\notin\overline\mu_2(y_{p-1}).                 \tag{20}
\]
The two theta witnesses and last-edge value follow from the swap and
prefix equality; theta's stored-set exclusion is (12) transported through
the two same-good relations; and epsilon is present at the penultimate
vertex because its earlier stage witness is distinct in the elementary
prefix and the component at \(y_{p-1}\) is disjoint from the changed
component at \(y_p\).  Thus (20) forms an actual
\noderef{HierarchySituationTwoCaseTwoData} for the final presentation,
contradicting \noderef{HierarchySituationTwoSubcaseTwoOne}.

The alternatives theta stored or nonstored, followed in the latter case
by alpha stored or nonstored, are exhaustive.  Every resulting branch has
reduced to the already established Subcase 1.1 or 2.1 contradiction, so
Subcase 3.2 is impossible.
\end{proof}
